\begin{abstract}
	Данная выпускная квалификационная работа посвящена разработке пайплайна для создания интерактивных цифровых двойников на основе воксельных медицинских данных. Традиционные методы визуализации КТ и МРТ, такие как прямой объемный рендеринг, требуют значительных вычислительных ресурсов для работы в реальном времени. В качестве альтернативы в работе исследуется применимость современной технологии 3D Gaussian Splatting (3DGS), которая позволяет совместить высокую скорость отрисовки с фотореалистичным качеством, что делает ее перспективным решением для прикладных медицинских и образовательных задач.
	
	\textbf{В целях исследования} был спроектирован процесс обработки данных на примере костной системы человека из анатомического набора данных Visible Human Project (сегментированного и очищенного другими исследователями). Разработанные скрипты преобразуют исходную сетку вокселей в формат OpenVDB для последующей генерации выборки ракурсов в графическом редакторе Blender. Интеграция процесса обучения с библиотекой gsplat обеспечивается за счет автоматического формирования конфигурационных файлов виртуальных камер в формате COLMAP, затем обученные трехмерные гауссианы сегментируются по принципу ближайшего соседа на основе исходной размеченной воксельной сетки.
	
	\textbf{Практическим результатом работы} стала фотореалистичная модель скелета человека и адаптированный под нее интерактивный браузерный просмотрщик на базе WebGL, где реализована функция динамического выделения цветом выбранного класса.
\end{abstract}